\begin{abstract}
The advent of research like BitNet~\cite{bitnet} is ushering in a transformative phase characterized by 1-bit Large Language Models (LLMs). In this study, we propose a novel 1-bit LLM variant, termed \textbf{\text{BitNet b1.58}}, where each parameter (or weight) within the model assumes ternary values from the set \{-1, 0, 1\}. This model attains parity with traditional full-precision Transformer LLMs (utilizing FP16 or BF16 precision) of equivalent architecture and trained on an identical volume of tokens, achieving similar perplexity and downstream task results. Moreover, it achieves substantial gains in efficiency by reducing latency, memory usage, throughput, and energy demands. Importantly, the introduction of the 1.58-bit LLM establishes an innovative scaling law and a novel framework for constructing future LLMs that retain high performance while enhancing cost-efficiency. Additionally, it paves the way for alternative computational strategies and sets the foundation for hardware designed explicitly to exploit the advantages of 1-bit LLMs.
\end{abstract}

\section{The Era of 1-bit LLMs}

Artificial Intelligence has recently witnessed unprecedented advancements with the escalating scale and abilities of Large Language Models (LLMs). While these models have delivered state-of-the-art outcomes across diverse natural language processing benchmarks, their growing size complicates real-world deployment and intensifies environmental and financial concerns owing to substantial energy requirements.

One promising direction to mitigate such challenges is post-training quantization, whereby pretrained models are converted into low-bit formats to facilitate inference~\cite{smoothquant, gptq, quip, quip_sharp}. By decreasing the numeric precision of both parameters and activations, memory and compute demands can be sharply diminished for LLMs. The field has progressively transitioned from 16-bit to lower-precision implementations, including prominent 4-bit schemes~\cite{gptq, awq}. Despite their broad adoption in industrial settings, however, post-training quantization remains sub-optimal relative to other approaches.

Recent advances in 1-bit neural architectures, exemplified by BitNet~\cite{bitnet}, indicate a compelling pathway for lowering LLM operating expenses without sacrificing accuracy. Conventional LLMs operate using 16-bit floating-point formats (FP16 or BF16), with the dominant computational demand arising from large-scale matrix multiplications, which inherently depend on floating-point arithmetic. BitNet, on the other hand, restricts its matrix computations to integer additions, resulting in dramatic reductions in energy consumption. Since chip performance is fundamentally constrained by power usage, this decrease in energy requirements also translates to an increase in computational speed.

Beyond computation, moving model parameters from off-chip DRAM into on-chip memory such as SRAM during inference introduces substantial costs. While expanding on-chip SRAM has been explored to improve throughput, the expense of SRAM compared to DRAM is significantly greater. 1-bit LLMs, however, boast a considerably smaller memory footprint than their full-precision counterparts, both in terms of storage requirements and memory bandwidth. This reduced size accelerates the transfer of weights from DRAM, streamlining and optimizing inference processes.

In this paper, we present a notable 1-bit LLM variant named \textbf{\text{BitNet b1.58}}, characterized by all model parameters being ternary, i.e., restricted to the values \{-1, 0, 1\}. By augmenting the standard 1-bit BitNet parameter set with an explicit 0, we attain an effective bit-width of 1.58 bits under a binary scheme. \text{BitNet b1.58} retains every advantage of its 1-bit predecessor, preserving the novel arithmetic framework that eliminates almost all multiplications in matrix operations and allowing for aggressive hardware-level optimizations. It sustains comparable energy and memory efficiency, as well as throughput and latency benefits, relative to traditional FP16-based LLMs. Additionally, \text{BitNet b1.58} introduces two unique improvements: first, its capacity for modeling is enhanced owing to explicit support for feature selection, as the presence of zeros in the weights facilitates more robust filtering—potentially yielding significant gains over traditional 1-bit LLMs. Second, empirical evaluation indicates that \text{BitNet b1.58} can achieve performance parity with full-precision (FP16) models in perplexity and downstream tasks for model sizes starting at 3B, given identical configurations (such as model size and number of training tokens).

\section{BitNet b1.58}

\text{BitNet b1.58} builds upon the BitNet architecture, a variant of the Transformer where the traditional \emph{nn.Linear} layers are replaced by \emph{BitLinear} layers. This model is trained from the beginning using weights quantized to 1.58 bits and activations quantized to 8 bits. Several enhancements differentiate it from the initial BitNet version, as outlined below.

\paragraph{Quantization Function.} The approach for weight quantization limits them to the values -1, 0, or +1 through the use of an \emph{absmean} quantization function. Here, the weight matrix is normalized by its mean absolute value prior to rounding each element to the nearest value within the set \{-1, 0, +1\}:

\begin{equation}
    \widetilde{W} = \mathrm{RoundClip}(\frac{W}{\gamma + \epsilon}, -1, 1),
\end{equation}

\begin{equation}
    \mathrm{RoundClip}(x, a, b) = \max(a, \min(b, \mathrm{round}(x))),
\end{equation}

\begin{equation}
    \gamma = \frac{1}{nm}\sum_{ij} |W_{ij}|.
\end{equation}

The procedure for activation quantization matches the strategy employed in BitNet, except that it omits scaling activations to $[0, Q_b]$ prior to applying non-linearities. Instead, activations for each token are mapped to $[-Q_b, Q_b]$, thereby eliminating the use of zero-point quantization. This modification not only streamlines both software implementation and hardware optimization but, based on our experiments, also incurs only negligible performance impact.

\paragraph{LLaMA-alike Components.}

The LLaMA~\cite{llama, llama2} design has become a foundational structure for the majority of open-source LLMs. To maximize compatibility with the broader community, \text{BitNet b1.58} incorporates components from LLaMA, namely RMSNorm~\cite{rmsnorm}, SwiGLU~\cite{swiglu}, and rotary embedding~\cite{rope}, while also discarding all bias terms. This alignment ensures that \text{BitNet b1.58} is readily deployable within widely used open-source ecosystems, such as Huggingface, vLLM~\cite{vllm}, and llama.cpp\footnote{\url{https://github.com/ggerganov/llama.cpp}}, requiring only minimal adaptation.

\section{Results}

We conducted a comparative analysis between \text{BitNet b1.58} and our reproduced FP16 \text{LLaMA LLM} models across multiple parameter scales. For consistency in evaluation, both models were pre-trained on the RedPajama dataset~\cite{redpajama}, each receiving an allocation of 100 billion tokens. The zero-shot capabilities of these models were assessed on a suite of language benchmarks, such as ARC-Easy~\cite{arc}, ARC-Challenge~\cite{arc}, Hellaswag~\cite{hellaswag}, Winogrande~\cite{winoGrande}, PIQA~\cite{piqa}, OpenbookQA~\cite{openbookqa}, and BoolQ~\cite{boolq}. Furthermore, we reported validation perplexity on the WikiText2~\cite{wikitext2} and C4~\cite{c4} corpora.

To evaluate efficiency, we measured GPU memory usage and inference latency for both \text{LLaMA LLM} and \text{BitNet b1.58}. These assessments utilized the FasterTransformer\footnote{\url{https://github.com/NVIDIA/FasterTransformer}} framework, designed to maximize inference throughput on GPU hardware. The 2-bit inference kernel from Ladder~\cite{ladder} was employed for \text{BitNet b1.58}. We present the average time required per generated token, which is the primary bottleneck during inference.

In Table~\ref{tab:ppl}, we present a summary of perplexity values and corresponding efficiency metrics for \text{BitNet b1.58} versus \text{LLaMA LLM}. The results indicate that \text{BitNet b1.58} achieves comparable perplexity to the full-precision \text{LLaMA LLM} at the 3B scale, while operating 2.71 times faster and utilizing 3.55 times less GPU memory. Notably, the 3.9B variant of \text{BitNet b1.58} offers a 2.4$\times$ speedup, consumes 3.32$\times$ less memory, and outperforms the 3B \text{LLaMA LLM} in accuracy.

Detailed zero-shot evaluation outcomes on downstream tasks are found in Table~\ref{tab:zero-shot}. The evaluation protocol was aligned with the procedure provided by \emph{lm-evaluation-harness}\footnote{\url{https://github.com/EleutherAI/lm-evaluation-harness}}. Results demonstrate that as model scale increases, the performance discrepancy between \text{BitNet b1.58} and \text{LLaMA LLM} diminishes. More critically, \text{BitNet b1.58} matches the accuracy of the full-precision baseline from 3B parameters onward. Consistent with perplexity trends, the 3.9B configuration of \text{BitNet b1.58} surpasses the 3B \text{LLaMA LLM} in end-task accuracy while operating with reduced memory and latency. These findings establish \text{BitNet b1.58} as a Pareto-optimal enhancement over contemporary LLM baselines.

\paragraph{Memory and Latency}

We extended our experiments by increasing the model sizes to 7B, 13B, and 70B parameters, and analyzed the associated costs. The patterns depicted in Figure~\ref{fig:latency-memory-energy} demonstrate that both latency and memory scale with model size, with acceleration benefits becoming more pronounced at larger scales. Specifically, \text{BitNet b1.58} 70B achieves a 4.1$\times$ speed advantage over the \text{LLaMA LLM} reference. This improvement arises because the computational time for \emph{nn.Linear} operations increases with larger models. Similarly, memory usage exhibits analogous scaling; as the embedding layer retains full-precision representation, its share of the overall memory footprint becomes less significant in larger architectures. Measurements for both latency and memory utilized a 2-bit kernel, suggesting that there is potential for further efficiency gains through additional optimizations.

\paragraph{Energy}

We also quantified the arithmetic operation energy demands for \text{BitNet b1.58} compared to \text{LLaMA LLM}. Our primary emphasis was on matrix multiplication calculations, which represent the dominant expense in large language models. The breakdown presented in Figure~\ref{fig:energy} reveals that \text{BitNet b1.58} relies primarily on INT8 addition, whereas \text{LLaMA LLM} leverages both FP16 addition and multiplication. Leveraging the energy consumption model from~\cite{energycost, pokebnn}, \text{BitNet b1.58} yields a 71.4-fold reduction in energy for matrix multiplications on 7nm hardware. We additionally provide end-to-end energy usage statistics for models processing 512 tokens. Our findings indicate that, with increasing model scale, \text{BitNet b1.58} demonstrates heightened energy efficiency relative to the FP16 \text{LLaMA LLM} baseline. This efficiency arises as the proportion of cost attributable to \emph{nn.Linear} operations grows with model size, whereas other parts of the model contribute less to the total energy requirement in larger networks.

\paragraph{Throughput}

We assessed the throughput of \text{BitNet b1.58} and the 70B-parameter \text{LLaMA LLM} on a setup consisting of two 80GB A100 GPUs, applying pipeline parallelism~\cite{gpipe} to enable \text{LLaMA LLM} 70B to fit on the hardware. The batch size was incrementally increased until the memory ceiling of the GPUs was reached, with each input having a sequence length of 512. According to Table~\ref{tab:throughput}, \text{BitNet b1.58} 70B supports batch sizes that are up to 11 times larger than those of \text{LLaMA LLM}, resulting in a throughput that is 8.9 times greater.

\textbf{\text{BitNet b1.58} paves the way for a novel scaling relationship between model capability and the cost of inference.} Drawing from Figures~\ref{fig:latency-memory-energy} and \ref{fig:energy}, the following correspondences illustrate equivalent efficiencies between models quantized to 1.58 bits and traditional 16-bit precision:

\begin{itemize}
    \item The 13B \text{BitNet b1.58} surpasses the 3B FP16 LLM in latency, memory footprint, and energy requirements.
    \item The 30B \text{BitNet b1.58} demonstrates better efficiency in latency, memory usage, and energy consumption compared to the 7B FP16 LLM.
    \item The 70B \text{BitNet b1.58} achieves higher efficiency across latency, memory, and energy when contrasted with the 13B FP16 LLM.
\end{itemize}

\paragraph{Training with 2T Tokens}

The amount of training data, measured in tokens, plays a key role in the effectiveness of LLMs. To investigate the token scalability of \text{BitNet b1.58}, we trained a version of this model on 2T tokens following the data preparation methodology of StableLM-3B~\cite{StableLM-3B-4E1T}, currently a leading open-source model at the 3B scale. Both models were benchmarked using a suite comprising Winogrande~\cite{winoGrande}, PIQA~\cite{piqa}, SciQ~\cite{sciq}, LAMBADA~\cite{lambada}, and ARC-easy~\cite{arc}. Zero-shot accuracy results are presented in Table~\ref{tab:2t}. For evaluations involving both accuracy and normalized accuracy, the mean value was reported. Performance figures for StableLM 3b at the 2T-token scale are reproduced from its original technical documentation. The evaluation reveals that \text{BitNet b1.58} consistently achieves higher scores across all tested benchmarks, suggesting that 1.58-bit LLMs not only maintain, but excel in, generalization performance.

\section{Discussion and Future Work}

\textbf{1-bit Mixture-of-Experts (MoE) LLMs}

Mixture-of-Experts (MoE) architectures have demonstrated notable efficiency gains for LLMs by lowering computation-related FLOPs. However, large memory usage and significant inter-chip communication requirements remain bottlenecks in practical applications. Utilizing 1.58-bit LLMs offers a promising solution to these issues. The lower memory demands translate directly to fewer devices needed for deploying MoE models and also lessen the communication costs associated with transmitting activations. In the ideal scenario, when the entire MoE model fits onto a single chip, inter-device communication becomes negligible.

\textbf{Native Support of Long Sequence in LLMs}

The growing need for LLMs to process extended sequences has introduced new challenges, with KV cache memory usage during inference being a primary concern. \text{BitNet b1.58} marks substantial progress toward addressing this issue by shrinking activation sizes from 16 bits down to 8 bits. This effectively doubles the manageable context length under identical hardware constraints. There exists further opportunity for improvement, such as compressing activation size to 4 bits or even fewer in the context of 1.58-bit LLMs, a direction we propose to pursue in future research.

\textbf{LLMs on Edge and Mobile}

Deploying 1.58-bit LLMs could substantially boost language model performance on resource-constrained edge and mobile hardware, where memory and computational throughput typically present limiting factors. The substantially lower memory and energy consumption profiles of 1.58-bit LLMs render them well suited for operation on these platforms, thus unlocking new applications not previously feasible on mobile or edge devices. These advances significantly broaden the functional scope and potential of edge computing. Additionally, as most edge and mobile systems rely on CPUs, the compatibility of \text{BitNet b1.58} with such processors permits efficient execution and further amplifies the capabilities of these devices.

\textbf{New Hardware for 1-bit LLMs}

Innovations such as those demonstrated by Groq\footnote{\url{https://groq.com/}} illustrate the viability and advantages of dedicated hardware solutions like LPUs for LLM deployment. Building on this foundation, we advocate for the development of new systems and hardware platforms specifically designed to capitalize on the unique computation characteristics of 1-bit LLMs, as enabled by the BitNet paradigm~\cite{bitnet}.